% TOP_PROBLEM: {"id": 30007426, "title": "Avoiding every dyadic affine sequence in a large-measure set", "category_id": 8, "category": {"id": 8, "name": "analysis", "display_name": "Analysis", "description": "Limits, continuity, calculus, and function theory.", "slug": "analysis", "order_index": 8, "created_at": "2026-07-31T15:26:25.671Z"}, "status": "open"}
% FIRST_POSTED: 2026-10-06
% ATTEMPT_STATUS: solved
% OPENAI_CLAIM: {"schema_version":1,"source_repo":"https://github.com/openai/math","source_commit":"adc7f1241b42e322a6451854ab7e4b4c146bf78a","release_date":"2026-10-06","match":"full","resolution":"proved","independently_reviewed":false,"second_pass":"agreed","families":[{"family":"084","title":"The geometric case of the Erdős similarity conjecture","subject":"Real and complex analysis","manuscripts":[{"title":"The dyadic case of the Erdős similarity conjecture","pdf_path":"preprints/The-dyadic-case-of-the-Erdos-similarity-conjecture-September-25-2026/paper.pdf","readme_path":"preprints/The-dyadic-case-of-the-Erdos-similarity-conjecture-September-25-2026/README.md","date":"2026-09-25","theorem_ref":"preprints/The-dyadic-case-of-the-Erdos-similarity-conjecture-September-25-2026/build/sections/introduction.tex :: thm:main"}],"lean":{"doc_path":"lean/docs/084.md","comparators":["lean/ComparatorChallenges/DyadicAvoidance.json","lean/ComparatorChallenges/DyadicAvoidance.lean"],"declarations":[],"covers_main_theorem":false},"reasoning_trace":null}]}
\documentclass[11pt]{article}
\usepackage[margin=25mm]{geometry}
\usepackage{fontspec}
\setmainfont{Latin Modern Roman}
\usepackage{amsmath,amssymb,mathtools}
\usepackage{unicode-math}
\setmathfont{Latin Modern Math}
\usepackage{enumitem}
\usepackage{xurl,hyperref,bookmark}
\hypersetup{colorlinks=true,urlcolor=blue}
\setcounter{secnumdepth}{0}
\setlength{\parindent}{0pt}
\setlength{\parskip}{5pt}
\setlength{\emergencystretch}{3em}
\begin{document}
\section*{Avoiding every dyadic affine sequence in a large-measure set}

\subsection{OpenAI's claim}
\textit{Avoiding affine dyadic sequences}: The dyadic paper claims compact subsets of $[0,1]$ with measure arbitrarily near one that omit at least one point of every translated, nontrivially dilated dyadic sequence. Positive and negative dilation parameters are both included.

\subsection{Outline of the argument}
\textit{Avoiding affine dyadic sequences}: A sparse periodic hitting set is built from random tables on an ordered finite tree. Routes preserved by suitable dyadic shifts provide many fresh terminal tests, while the edge-wise scale count is beaten by their failure probability. A small open repair handles exceptional centers; rescaling and reflection extend the hitting property to all nonzero signed scales.

\subsection{Scope relative to this problem}
\textit{Avoiding affine dyadic sequences}: The primary statement says: For every eta>0 there is a compact subset of [0,1] of measure greater than 1-eta containing no affine copy of \{2\textasciicircum{}(-n):n>=1\}, with either sign of nonzero dilation. The target retains these quantified objects and hypotheses, so its positive assertion is covered.

With the hypotheses and conventions above, the claimed result addresses the whole mathematical assertion on this page.

\subsection{Formal verification}
Family 084 describes Lean coverage for avoidance for the dyadic ratio; it does not assert avoidance for an arbitrary geometric ratio.

\textit{Avoiding affine dyadic sequences}: The manuscript links \texttt{DyadicAvoidance.lean}.

The documented selection does not cover the entire manuscript assertion relevant here.

This site has not run the Lean development. These links describe OpenAI’s selected formal statements and should be read with their stated scope.

\subsection{Status caveat}
OpenAI reports that an internal model produced the manuscripts. Their arguments have not been independently verified as proofs by this site. Any portion outside the documented formal scope may contain mistakes, and the catalogue label reports the announced mathematical claim.

\subsection{Sources}
\begin{itemize}
\item OpenAI, \emph{The dyadic case of the Erdős similarity conjecture}, 2026-09-25; citation key \texttt{OAI:The-dyadic-case-of-the-Erdos-similarity-conjecture-September-25-2026}. \href{https://github.com/openai/math/blob/adc7f1241b42e322a6451854ab7e4b4c146bf78a/preprints/The-dyadic-case-of-the-Erdos-similarity-conjecture-September-25-2026/paper.pdf}{PDF (GitHub)}; \href{https://github.com/openai/math/raw/adc7f1241b42e322a6451854ab7e4b4c146bf78a/preprints/The-dyadic-case-of-the-Erdos-similarity-conjecture-September-25-2026/paper.pdf}{PDF (direct)}; \href{https://github.com/openai/math/blob/adc7f1241b42e322a6451854ab7e4b4c146bf78a/preprints/The-dyadic-case-of-the-Erdos-similarity-conjecture-September-25-2026/README.md}{Manuscript page and citation}; \href{https://github.com/openai/math/blob/main/preprints/The-dyadic-case-of-the-Erdos-similarity-conjecture-September-25-2026/paper.pdf}{Latest version}.
\item Family 084: \href{https://github.com/openai/math/blob/adc7f1241b42e322a6451854ab7e4b4c146bf78a/lean/docs/084.md}{Lean scope}.
\item \href{https://github.com/openai/math/blob/adc7f1241b42e322a6451854ab7e4b4c146bf78a/lean/ComparatorChallenges/DyadicAvoidance.json}{Comparator DyadicAvoidance.json}.
\item \href{https://github.com/openai/math/blob/adc7f1241b42e322a6451854ab7e4b4c146bf78a/lean/ComparatorChallenges/DyadicAvoidance.lean}{Comparator DyadicAvoidance.lean}.
\item \href{https://github.com/openai/math/blob/adc7f1241b42e322a6451854ab7e4b4c146bf78a/CONTENTS.md}{OpenAI catalogue entry}.
\end{itemize}
\end{document}
